% BHU PYQ -- Professional Exam Template
\documentclass[11pt,a4paper]{article}

\usepackage[a4paper,top=2.5cm,bottom=2.5cm,left=2.8cm,right=2.8cm,headheight=14pt]{geometry}
\usepackage{amsmath,amssymb}
\usepackage[T1]{fontenc}
\usepackage[utf8]{inputenc}
\usepackage{lmodern}
\usepackage{microtype}
\usepackage{fancyhdr}
\usepackage{enumitem}
\usepackage{listings}
\usepackage{hyperref}
\usepackage{lastpage}
\usepackage{parskip}

\hypersetup{colorlinks=true,urlcolor=black,linkcolor=black,
  pdftitle={BVCA-404: Microprocessor -- BHU PYQ}}

\pagestyle{fancy}
\fancyhf{}
\renewcommand{\headrulewidth}{0.4pt}
\renewcommand{\footrulewidth}{0.4pt}
\fancyhead[L]{\small\textit{Banaras Hindu University}}
\fancyhead[R]{\small\textit{BVCA-404: Microprocessor}}
\fancyfoot[C]{\small Page \thepage\ of \pageref{LastPage}}

\lstset{
  basicstyle=\small\ttfamily,
  frame=single,
  breaklines=true,
  columns=flexible,
  keepspaces=true
}

\newlist{parts}{enumerate}{2}
\setlist[parts,1]{label=(\alph*),leftmargin=*,itemsep=4pt,topsep=3pt}
\setlist[parts,2]{label=(\roman*),leftmargin=*,itemsep=2pt,topsep=2pt}
\newcommand{\pts}[1]{\hfill\textit{[#1]}}

\begin{document}
\begin{center}
  {\large\bfseries BANARAS HINDU UNIVERSITY}\\[2pt]
  {\normalsize B.Voc. Semester IV Examination, 2023-24}\\[6pt]
  \rule{0.8\linewidth}{0.5pt}\\[6pt]
  {\large\bfseries Paper: BVCA-404 --- Microprocessor}\\[4pt]
  {\normalsize Time: Three Hours \hfill Maximum Marks: 70}
\end{center}

\rule{\linewidth}{0.4pt}

\smallskip
\noindent\textit{Instructions: Answer the questions in each section as instructed. All questions carry marks as indicated.}

\bigskip

\noindent{\large\bfseries SECTION --- A}
\rule{\linewidth}{0.4pt}\smallskip

Long-Answer Questions. Answer any two questions: \hfill \textbf{2 $\times$ 12 = 24 Marks}

\medskip
\noindent\textbf{Question 1.} Explain the pin diagram of the 8085 microprocessor in detail.

\medskip
\noindent\textbf{Question 2.} Describe the internal block diagram of the 8085 microprocessor.

\medskip
\noindent\textbf{Question 3.} What are interrupts? Differentiate between software interrupts and hardware interrupts in the 8085 microprocessor.

\medskip
\noindent\textbf{Question 4.} How many types of addressing modes does the 8085 microprocessor have? Explain any two of them with suitable examples.

\bigskip\noindent{\large\bfseries SECTION --- B}
\rule{\linewidth}{0.4pt}\smallskip

\noindent\textbf{Question 5.} Attempt any six parts (Answer in about 200 words each): \hfill \textbf{6 $\times$ 6 = 36 Marks}
\begin{parts}
  \item Explain the three-step instruction execution cycle (Fetch, Decode, Execute) in detail.
  \item What are universal logic gates? Explain why they are called so.
  \item Differentiate between a microprocessor, a microcontroller, and a microcomputer.
  \item What is assembly language? Discuss its advantages over machine language.
  \item Explain the computer memory hierarchy diagram in detail.
  \item Explain all three system buses: address bus, data bus, and control bus.
  \item Describe the instruction format and instruction types of the 8085 microprocessor.
  \item Draw a block diagram representing a microprocessor-based system and explain its basic components and functionality.
  \item Explain the architectural block diagram of the 8085 microprocessor.
\end{parts}

\bigskip\noindent{\large\bfseries SECTION --- C}
\rule{\linewidth}{0.4pt}\smallskip

\noindent\textbf{Question 6.} Attempt all parts: \hfill \textbf{10 $\times$ 1 = 10 Marks}
\begin{parts}
  \item What is a microprocessor?
  \item Write the list of logical operations performed by the 8085 microprocessor.
  \item What are the three main units that a microprocessor-based system is made of?
  \item Convert binary number ``101101110'' into its decimal equivalent.
  \item Draw the logic gate symbol that represents a NAND gate.
  \item Find out the 2's complement of the binary number ``001001110''.
  \item What is machine language?
  \item Is the pin diagram of the 8085 microprocessor considered an external diagram or an internal diagram?
  \item Write down the functionality of the $\overline{\text{RD}}$ and $\overline{\text{WR}}$ signals in brief.
  \item What are maskable interrupts in the 8085 microprocessor?
\end{parts}

\vfill
\rule{\linewidth}{0.4pt}
\begin{center}\small
\href{https://bhu-exam.vercel.app/}{Click here to visit our website}\\[4pt]\href{https://me-aryan.vercel.app/}{Click here to visit our contributor}\\[4pt]\href{https://quashan-me.vercel.app/}{Click here to visit our contributor}
\end{center}
\end{document}
